\documentclass{article}

%Basics
\usepackage[margin=1in]{geometry}
\usepackage{amsmath, amssymb, amsthm}
\usepackage{enumitem}

%Formatting and Spacing
\setitemize[1]{noitemsep, parsep = 5pt, topsep = 5pt}
\setenumerate[1]{label = (\alph*), parsep = 1pt, topsep = 5pt}
\setlength\parindent{0pt}
\linespread{1.1}

%Easy Numbering
\newcounter{counter}
\newcommand{\Exercise}{Exercise \stepcounter{counter}\thecounter}

%Cases Environment (Messy)
\newlist{Cases}{enumerate}{3}
\setlist[Cases]{leftmargin = .25in, label = {Case \arabic*.}, topsep = 0.01in, itemsep = 0.04in, itemindent = .5in, parsep = 0in}

%Custom Title Fields
\newcommand{\hmwkTitle}{Extra Homework (with Solutions)}
\newcommand{\hmwkDueDate}{February 31, 2000}
\newcommand{\hmwkClass}{Honors Discrete Mathematics}
\newcommand{\hmwkClassInstructor}{Gerandy Brito}
\newcommand{\hmwkSection}{Spring 2022}
\newcommand{\hmwkAuthorName}{Sarthak Mohanty}

%Headers and Footers
\usepackage{fancyhdr}
\usepackage{extramarks}
\pagestyle{fancy}
\lhead{\hmwkAuthorName}
\chead{\hmwkClass \ (\hmwkClassInstructor)}
\rhead{Homework 7}
\cfoot{\thepage}
\renewcommand\headrulewidth{0.4pt}
\renewcommand\footrulewidth{0.4pt}

\title{
    \vspace{2in}
    \textbf{\hmwkClass:\\ \hmwkTitle}\\
    \normalsize\vspace{0.1in}\small{Due\ on\ \hmwkDueDate\ at 11:59pm}\\
    \vspace{0.1in}\large{\textit{\hmwkClassInstructor\ \hmwkSection}}
    \vspace{3in}
    \author{\textbf{\hmwkAuthorName}}
    \date{}
}

\begin{document}

\maketitle
\pagebreak

\subsection*{\Exercise}
    For each part, show all your computations and explain how you used Fermat's theorem.
    \begin{enumerate}
        \item Compute $10^{33} \pmod{7}$.
        \item Show that $(x + y)^{p} \equiv x^{p} + y^{p} \pmod{p}$, where $p$ is prime.
        \item Compute $2^{p-2} + 3^{p-2} + 6^{p-2} - 1 \pmod{p}$, where $p > 3$ is prime.
    \end{enumerate}

\subsection*{Solution}
    \begin{enumerate}
        \item Since $7$ is prime, by Fermat's theorem, $10^{6} \equiv 1 \pmod{7}$. Hence $10^{33} = 10^{3}(10^{6})^{5} \equiv 10^{3} \equiv 6 \pmod{7}$.
        \item Since $p$ is prime, by Fermat's theorem, $(x + y)^{p} \equiv x + y \pmod{p}$, $x \equiv x^{p} \pmod{p}$ and $y \equiv y^{p} \pmod{p}$. Hence $(x + y)^{p} \equiv x + y \equiv x^{p} + y^{p} \pmod{p}$, and we have our desired result.
        
        \vspace{1.5mm}
        \textbf{Remark}: This congruency is commonly known as the Freshman's Dream, due to the erroneous calculation $(x + y)^{p} = x^{p} + y^{p}$ that many beginning students make.
        \item Since $p$ is prime, by Fermat's theorem, $2^{p - 1} \equiv 1 \pmod{p}$, $3^{p - 1} \equiv 1 \pmod{p}$, and $6^{p - 1} \equiv 1 \pmod{p}$. Then
        $$\begin{aligned}[t]
            &2^{p-2} + 3^{p-2} + 6^{p-2} - 1 \\
            ={}& (2^{p - 1} \cdot 2^{-1}) + (3^{p - 1} \cdot 3^{-1}) + (6^{p - 1} \cdot 6^{-1}) - 1 \\
            \equiv{}& 2^{-1} + 3^{-1} + 6^{-1} - 1 \pmod{p}. \\
        \end{aligned}$$
        Now $p > 2$, so $p \nmid 2$. As a result, $2 \mid p + 1$, so $\frac{p + 1}{2} \in \mathbb{Z}$, and thus $2^{-1} \equiv \frac{p + 1}{2} \pmod{p}$. Furthermore, since $p > 3$, $p \nmid 3$. Hence $3 \mid p + 1$ or $3 \mid p - 1$.
        
        \begin{Cases}
            \item Suppose $3 \mid p + 1$. Then $\frac{p + 1}{3} \in \mathbb{Z}$, and thus $3^{-1} \equiv \frac{p + 1}{3} \pmod{p}$. Furthermore, since $2 \mid p + 1$, $p + 1$ is divisible by $6$. Then $\frac{p + 1}{6} \in \mathbb{Z}$, and thus $6^{-1} \equiv \frac{p + 1}{6} \pmod{p}$. We conclude that $2^{-1} + 3^{-1} + 6^{-1} - 1 \equiv \frac{p + 1}{2} + \frac{p + 1}{3} + \frac{p + 1}{6} - 1 = 6p \equiv 0 \pmod{p}$.
            \item Now suppose $3 \mid p - 1$. Then $\frac{p - 1}{3} \in \mathbb{Z}$, and thus $3^{-1} \equiv \frac{p - 1}{3} \pmod{p}$. Furthermore, since $2 \mid p - 1$, $p - 1$ is divisible by $6$. Then $\frac{p - 1}{6} \in \mathbb{Z}$, and thus $6^{-1} \equiv \frac{p - 1}{6} \pmod{p}$. We conclude that $2^{-1} + 3^{-1} + 6^{-1} - 1 \equiv \frac{p + 1}{2} + \frac{p - 1}{3} + \frac{p - 1}{6} - 1 = 6p \equiv 0 \pmod{6}$.
        \end{Cases}
        In either case, $2^{p-2} + 3^{p-2} + 6^{p-2} - 1 \equiv 0$.
    \end{enumerate}

\subsection*{\Exercise}
    Use congruencies to prove each part below. Show your work.
    \begin{enumerate}
        \item Show that $5$ divides $6^{n} - 1$ for all natural numbers $n$.
        \item Show that $7$ divides $3^{n} + 4^{n}$ for all odd numbers $n$.
        \item Show that $7$ divides $111\dots1$, where there are 2022 many ones.
    \end{enumerate}

\subsection*{Solution}
    \begin{enumerate}
        \item Note that $6 \equiv 1 \pmod{5}$, so $6^{n} \equiv 1 \pmod{5}$, and by definition $5 \mid 6^{n} - 1$ for all $n \in \mathbb{N}$.
        \item Note that $3 \equiv -4 \pmod{7}$, so $3^{n} \equiv (-4)^{n} = (-1)^{n}(4^{n}) \pmod{7}$. Since $n$ is odd, $(-1)^{n} = -1$, so $3^{n} \equiv -4^{n} = 0 \pmod{7}$, so $3^{n} + 4^{n} \equiv 0 \pmod{7}$, and by definition $7 \mid 3^{n} + 4^{n}$.
        \item Let $n$ be the number $$n = 111\dots 1\text{, where there are 2022 many ones}.$$ Then 
        \begin{equation*}
            n = \sum_{i = 0}^{2021} 10^{i} = \sum_{i = 0}^{\left\lfloor \frac{2021}{6} \right\rfloor} \sum_{j = 0}^{5} 10^{6i + j} = \sum_{i = 0}^{336} (111111 \times 10^{6i}).
        \end{equation*}
        
        Now since $7$ is prime, by Fermat's Theorem, $10^{6} \equiv 1 \pmod{7}$. Furthermore, $111111 \equiv 0 \pmod{7}$. Hence $111111 \times 10^{6} \equiv 1 \cdot 0 = 0 \pmod{7}$, so $$n = \sum_{i = 0}^{336} (111111 \times 10^{6i}) = \sum_{i = 0}^{336} (111111 \times 10^{6})^{i} \equiv 0 \pmod{7},$$ and by definition $7 \mid n$.
    \end{enumerate}

\subsection*{\Exercise}
    \textbf{Wilson's theorem} says that a number $p$ is prime if and only if $$(p - 1)! + 1 \equiv 0 \pmod{p}.$$
    \begin{enumerate}
        \item Recall that for $p$ prime, every number $0 < x < p$ has an inverse (mod $p$). Which of these numbers are their own inverse? Prove your answer.
        \item Use part (a) to simplify $(p - 1)! \pmod{p}$ and conclude that, for $p$ prime, $(p - 1)! + 1 \equiv 0 \pmod{p}$. Show your work.
        \item Conclude the proof of the theorem by showing that, if $N$ is not prime, then $(N - 1)! + 1$ is not a multiple of $N$. (\textit{Hint: since $N$ is not prime, there is a divisor $d$ of $N$ with $1 < d < N$}.)
    \end{enumerate}

\subsection*{Solution}
    \begin{enumerate}
        \item If $x$ is its own inverse, then $x^{2} \equiv 1 \pmod{p}$, so $x^{2} - 1 = (x + 1)(x - 1) \equiv 0 \pmod{p}$, so $x = 1$ or $x = p - 1$.
        \item Since $p$ is prime, each $a \in \{1, \dots p - 1\}$ has a unique inverse $a^{-1} \in \{1, \dots p - 1\}$. As shown in part (a), $a = a^{-1}$ for $a = 1$ and $a = p - 1$. Now we rearrange the factorization of $(p - 1)!$ so that each $a \in \{2, \dots, p - 2\}$ is matched to its unique inverse modulo $p$, leading to
        $$(p - 1)! = 1 \times 2 \times \dots \times (p - 2) \times (p - 1) \equiv 1 \times 1 \times (p - 1) \equiv p - 1 \pmod{p}.$$ Hence $(p - 1)! + 1 \equiv (p - 1) + 1 = p \equiv 0 \pmod{p}$.
        \item Suppose $N$ is composite. Then there exists some $d$ such that $d \mid N$ and $1 < d < N$. Now since $d \mid N$ and $d \ne N$, it follows that $d \mid (N - 1)!$. 
        
        \vspace{1.5mm}
        We will conclude this proof using contradiction. Suppose $(N - 1)! + 1$ is a multiple of $N$. Then $kN = (N - 1)! + 1$, or equivalently, $kN - (N - 1)! = 1$ for some integer $k$. But $d \mid N$ and $d \mid (N - 1)!$, so $d \mid 1$. This is a contradiction, as $1 < d < N$. Hence $(N - 1)! + 1$ is not a multiple of $N$.
    \end{enumerate}

\subsection*{\Exercise}
    (Application of Wilson's Theorem)
    \begin{enumerate}
        \item Let $$1 + \frac{1}{2} + \frac{1}{3} + \dots + \frac{1}{23} = \frac{a}{23!}.$$ Find the remainder when $a$ is divided by $13$.
        \item Find $$\gcd(19! + 19, 20! + 19).$$
        \item Let $p$ be an odd prime. Prove that
        $$\left(\left(\frac{p - 1}{2}\right)!\right)^2 \equiv (-1)^{\frac{p + 1}{2}} \pmod{p}.$$
    \end{enumerate}
    
\subsection*{Solution}
    \begin{enumerate}
        \item Multiplying both sides by $23!$, we get
        $$a = \frac{23!}{1} + \frac{23!}{2} + \dots + \frac{23!}{23}.$$
        We want to compute $a \pmod{13}$. Note that every term except $\frac{23!}{13}$ still has a factor of $13$, so all those terms are $0 \pmod{13}$. Hence
        $$a \equiv \frac{23!}{13} \pmod{13}.$$
        Now
        $$\frac{23!}{13} = 1 \cdot 2 \cdots 12 \cdot 14 \cdot 15 \cdots 23.$$
        Reducing modulo $13$, we get
        $$\frac{23!}{13} \equiv 12! \cdot 10! \pmod{13}.$$
        By Wilson's Theorem,
        $$12! \equiv -1 \pmod{13}.$$
        Also,
        $$10! \cdot 11 \cdot 12 = 12! \equiv -1 \pmod{13}.$$
        Since $11 \cdot 12 \equiv (-2)(-1) \equiv 2 \pmod{13}$, we have
        $$2 \cdot 10! \equiv -1 \pmod{13}.$$
        Since $2^{-1} \equiv 7 \pmod{13}$,
        $$10! \equiv -7 \pmod{13}.$$
        Therefore,
        $$a \equiv 12! \cdot 10! \equiv (-1)(-7) \equiv 7 \pmod{13}.$$
        Hence the remainder is
        $$\boxed{7}.$$

        \item Let
        $$A = 19! + 19, \qquad B = 20! + 19.$$
        Then
        $$B - 20A = (20! + 19) - 20(19! + 19) = -361.$$
        Hence
        $$\gcd(A,B) \mid 361.$$
        Now
        $$A = 19! + 19 = 19(18! + 1).$$
        By Wilson's Theorem,
        $$18! \equiv -1 \pmod{19},$$
        so $19 \mid 18! + 1$. Thus
        $$19^2 \mid A.$$
        Also,
        $$B = 20! + 19 = 19(20 \cdot 18! + 1).$$
        Since $20 \equiv 1 \pmod{19}$ and $18! \equiv -1 \pmod{19}$,
        $$20 \cdot 18! + 1 \equiv -1 + 1 \equiv 0 \pmod{19}.$$
        Thus
        $$19^2 \mid B.$$
        Therefore $361 \mid \gcd(A,B)$, and since $\gcd(A,B) \mid 361$, we conclude that
        $$\boxed{\gcd(19! + 19, 20! + 19) = 361}.$$

        \item By Wilson's Theorem,
        $$(p - 1)! \equiv -1 \pmod{p}.$$
        Now write $(p - 1)!$ by pairing the first half of the factors with the second half:
        $$(p - 1)! = 1 \cdot 2 \cdots \frac{p - 1}{2} \cdot \left(\frac{p + 1}{2}\right) \cdots (p - 2)(p - 1).$$
        Modulo $p$, the second half is
        $$\frac{p + 1}{2}, \frac{p + 3}{2}, \dots, p - 1
        \equiv -\frac{p - 1}{2}, -\frac{p - 3}{2}, \dots, -1 \pmod{p}.$$
        Hence
        $$(p - 1)! \equiv \left(\frac{p - 1}{2}\right)! \cdot (-1)^{\frac{p - 1}{2}}\left(\frac{p - 1}{2}\right)! \pmod{p}.$$
        Therefore,
        $$(p - 1)! \equiv (-1)^{\frac{p - 1}å{2}}\left(\left(\frac{p - 1}{2}\right)!\right)^2 \pmod{p}.$$
        Using Wilson's Theorem again,
        $$(-1)^{\frac{p - 1}{2}}\left(\left(\frac{p - 1}{2}\right)!\right)^2 \equiv -1 \pmod{p}.$$
        Multiplying both sides by $(-1)^{\frac{p - 1}{2}}$, we obtain
        $$\left(\left(\frac{p - 1}{2}\right)!\right)^2 \equiv (-1)^{\frac{p + 1}{2}} \pmod{p}.$$
    \end{enumerate}


\subsection*{\Exercise}
    Find all pairs of natural numbers $(x, p)$ where $p$ is prime, such that $$x \text{ (mod }p) + x \text{ (mod }2p) + x \text{ (mod }3p) + x \text{ (mod }4p) = x + p.$$
    Here, $a \text{ (mod }b)$ denotes the remainder of the division of $a$ by $b$.

\subsection*{Solution}
    We write
    $$x = pq + r,$$
    where $0 \le r < p$. Then
    $$x \text{ (mod }p) = r.$$
    More generally, for $k = 1, 2, 3, 4$,
    $$x \text{ (mod }kp) = p(q \text{ (mod }k)) + r.$$
    Thus the given equation becomes
    $$r + \big(p(q \text{ (mod }2)) + r\big)
    + \big(p(q \text{ (mod }3)) + r\big)
    + \big(p(q \text{ (mod }4)) + r\big)
    = pq + r + p.$$
    Simplifying,
    $$3r = p\Big(q + 1 - (q \text{ (mod }2)) - (q \text{ (mod }3)) - (q \text{ (mod }4))\Big).$$

    First suppose $p \neq 3$. Since $p$ is prime, $p \nmid 3$. Hence $p \mid r$. But $0 \le r < p$, so $r = 0$. Therefore,
    $$q + 1 = (q \text{ (mod }2)) + (q \text{ (mod }3)) + (q \text{ (mod }4)).$$
    Since the right hand side is at most $1 + 2 + 3 = 6$, it suffices to check $q = 0, 1, \dots, 5$:
    $$
    \begin{array}{c|c|c}
        q & (q \text{ (mod }2)) + (q \text{ (mod }3)) + (q \text{ (mod }4)) & q + 1 \\ \hline
        0 & 0 & 1 \\
        1 & 3 & 2 \\
        2 & 4 & 3 \\
        3 & 4 & 4 \\
        4 & 1 & 5 \\
        5 & 4 & 6
    \end{array}
    $$
    Thus $q = 3$, and so $x = 3p$.

    Now suppose $p = 3$. Then
    $$r = q + 1 - (q \text{ (mod }2)) - (q \text{ (mod }3)) - (q \text{ (mod }4)).$$
    Since $0 \le r < 3$, it suffices to check when the right hand side is $0$, $1$, or $2$. Also, since the sum of the remainders is at most $6$, we only need to check $q \le 7$:
    $$
    \begin{array}{c|c|c}
        q & (q \text{ (mod }2)) + (q \text{ (mod }3)) + (q \text{ (mod }4)) & r \\ \hline
        0 & 0 & 1 \\
        1 & 3 & -1 \\
        2 & 4 & -1 \\
        3 & 4 & 0 \\
        4 & 1 & 4 \\
        5 & 4 & 2 \\
        6 & 2 & 5 \\
        7 & 4 & 4
    \end{array}
    $$
    Hence the possible pairs $(q,r)$ are $(0,1)$, $(3,0)$, and $(5,2)$. Since $x = 3q + r$, this gives $x = 1, 9, 17$.

    Therefore, the solutions are
    $$\boxed{(x,p) = (3p,p) \text{ for all primes } p \neq 3}$$
    together with
    $$\boxed{(x,p) = (1,3), (9,3), (17,3)}.$$

\subsection*{\Exercise}
    \textit{(RSA Implementation)} First express your answers without computing modular exponentiations. Then use a computational aid to complete these computations.
    \begin{enumerate}
        \item Encrypt the message ATTACK using the RSA system with $n = 47 \cdot 59$ and $e = 13$, translating each letter into integers and grouping together pairs of integers, as in Chapter 4.6 Example 8 in the textbook.
        \item What is the original message encrypted using the RSA system with $n = 53 \cdot 61$ and $e=17$ if the encrypted message is $3185\ 2038\ 2460\ 2550$? (To decrypt, first find the decryption exponent $d$.)
    \end{enumerate}

\subsection*{Solution}
    \begin{enumerate}
        \item First, we translate the letters in ATTACK into their numerical equivalents. Using $A = 00, B = 01, \dots, Z = 25$, we get
        $$\text{ATTACK} \mapsto 00 \quad 19 \quad 19 \quad 00 \quad 02 \quad 10.$$
        Since
        $$n = 47 \cdot 59 = 2773,$$
        and
        $$2525 < 2773 < 252525,$$
        we group the message into blocks of four digits:
        $$0019 \quad 1900 \quad 0210.$$
        We now encrypt each block using
        $$C \equiv M^{13} \pmod{2773}.$$
        Before computing, this gives
        $$19^{13} \pmod{2773}, \qquad 1900^{13} \pmod{2773}, \qquad 210^{13} \pmod{2773}.$$

        Since $2773 = 47 \cdot 59$, we compute modulo $47$ and modulo $59$, then use the Chinese Remainder Theorem. This gives
        $$
        \begin{array}{c|c|c|c}
            M & M^{13} \pmod{47} & M^{13} \pmod{59} & M^{13} \pmod{2773} \\ \hline
            19 & 35 & 57 & 411 \\
            1900 & 15 & 19 & 1848 \\
            210 & 31 & 52 & 642
        \end{array}
        $$
        Therefore, padding with zeroes when necessary, the encrypted message is
        $$\boxed{0411 \quad 1848 \quad 0642}.$$

        \item We first calculate $d$ as the inverse of $17 \pmod{52 \cdot 60}$. Since
        $$52 \cdot 60 = 3120,$$
        we want
        $$17d \equiv 1 \pmod{3120}.$$
        By the Euclidean algorithm,
        $$3120 = 17 \cdot 183 + 9, \qquad 17 = 9 \cdot 1 + 8, \qquad 9 = 8 \cdot 1 + 1.$$
        Reversing this,
        $$1 = 9 - 8 = 9 - (17 - 9) = 2 \cdot 9 - 17 = 2(3120 - 17 \cdot 183) - 17.$$
        Thus
        $$1 = 2 \cdot 3120 - 17 \cdot 367,$$
        so
        $$17^{-1} \equiv -367 \equiv 2753 \pmod{3120}.$$
        Hence $d = 2753$.

        We now decrypt each block using
        $$M \equiv C^{2753} \pmod{3233}.$$
        Since
        $$3233 = 53 \cdot 61,$$
        we reduce modulo $53$ and modulo $61$. By Fermat's theorem,
        $$C^{52} \equiv 1 \pmod{53}, \qquad C^{60} \equiv 1 \pmod{61}.$$
        Also,
        $$2753 \equiv 49 \pmod{52}, \qquad 2753 \equiv 53 \pmod{60}.$$
        Hence
        $$C^{2753} \equiv C^{49} \pmod{53}, \qquad C^{2753} \equiv C^{53} \pmod{61}.$$

        Using repeated squaring modulo $53$ and $61$, then recombining with the Chinese Remainder Theorem, we get
        $$
        \begin{array}{c|c|c|c}
            C & C^{49} \pmod{53} & C^{53} \pmod{61} & C^{2753} \pmod{3233} \\ \hline
            3185 & 14 & 47 & 1816 \\
            2038 & 47 & 56 & 2008 \\
            2460 & 21 & 9 & 1717 \\
            2550 & 40 & 45 & 411
        \end{array}
        $$
        Therefore, the decrypted numerical message is
        $$1816 \quad 2008 \quad 1717 \quad 0411.$$
        Translating this back to English letters,
        $$18 \quad 16 \quad 20 \quad 08 \quad 17 \quad 17 \quad 04 \quad 11
        \mapsto \text{SQUIRREL}.$$
        Hence the original message is
        $$\boxed{\text{SQUIRREL}}.$$
    \end{enumerate}

\subsection*{\Exercise} 
    (\textit{RSA Validity \& Uniqueness}) Let $(N, e)$ be a public key. In this problem, we will check that we can recover the correct message.
    \begin{enumerate}
        \item Explain why if the message $m$ is greater than $N$ we end up losing information.
        
        \vspace{3mm}
        One way to avoid the above situation is to break the message is to break the message into smaller parts. Let $x < N$ and $\gcd(x, N) = 1$. Assume there are two such messages, $x$ and $y$, such that after decoding $x^{e}$ and $y^{e}$ we get the same message.
        \item Check that, for $x, y$ as above, if $x^{e} \equiv y^{e} \pmod{N}$ then we must have $x = y$. 
        
        \vspace{1mm} \textit{Hint: because $\gcd(e, \phi(N)) = 1$ we can compute the inverse of $e, d$, mod $\phi(N)$. Consider $x^{ed}$ and $y^{ed}$}.)
    \end{enumerate}
    
    \textbf{Remark. }Although not graded, I strongly recommend you fully understand the proof of the decryption step of the RSA process, which can be found in Chapter 4.6 of the textbook.

\subsection*{Solution}
    \begin{enumerate}
        \item Consider $m \equiv c^{d} \pmod{N}$, the decryption formula used to obtain the original message. If $m > N$, then $c^{d} < N < m$ and thus we will obtain a number smaller than our original message $m$, causing a loss in information.
        \item Suppose $x^{e} \equiv y^{e} \pmod{N}$. Then $x^{ed} \equiv y^{ed} \pmod{N}$, so $x \equiv y \pmod{N}$. But $x, y < N$, so $x = y$.
    \end{enumerate}

\subsection*{\Exercise}
    \textit{(RSA Vulnerabilities)}
    \begin{enumerate}
        \item Explain why we should never use $(N = 2q, e)$ as a public key in an RSA system.
        \item Let $(N, e = 3)$ be your public key. Suppose you want to encrypt the message $m$ such that $m^{3} < N$. Explain why sending this small message is not protected by the RSA system.
        \item Prof.\ Brito's kids, Dario, Oliver, and Raquel, each use an RSA system with public keys $(N_{i}, e = 3)$, for $1 \le i \le 3$ to communicate with him. Brito sends the same message $m$ to each of them using their corresponding RSA system. That is, he sends $m^{3} \pmod{N_{i}}$ to each kid. Lianet (Brito's wife!), intercepted the three encrypted messages! Show how Lianet can find the original message $m$. You must explain a procedure Lianet can use to find $m$.
    \end{enumerate}

\subsection*{Solution}
    \begin{enumerate}
        \item One of the factors of $N = 2q$ is $2$. Due to the parity of $N$, it will be easy for an attacker to guess this number and therefore obtain the factorization of $N$, compromising the security of our RSA system.
        \item Suppose there exists some encrypted message $c$. Then $m^{3} \equiv c \pmod{N}$, but since $m^{3} < N$, $m^{3} = c$. Hence an attacker can easily obtain the original message, given by $m = \sqrt[3]{c}$.
        \item Let $c_{i}$ be the encrypted message sent to each child, and assume all $N_{i}$ are pairwise coprime. Then Lianet has three encrypted messages $c_{1}$, $c_{2}$, and $c_{3}$, where
        $$\begin{aligned}[t]
            c_{1} &\equiv m^{3} \pmod{N_{1}} \\
            c_{2} &\equiv m^{3} \pmod{N_{2}} \\
            c_{3} &\equiv m^{3} \pmod{N_{3}}.
        \end{aligned}$$
        Let $N = N_{1}N_{2}N_{3}$. By the Chinese Remainder Theorem, there exists some $c$ such that $$c \equiv m^{3} \pmod{N}$$ that is unique for all values less than $N$. But since $m$ is chosen such that $m < N_{1}, N_{2}, N_{3}$, it follows that $m^{3} < N$. With this knowledge, we conclude that $c = m^{3}$. Lianet can use the Extended Euclidean Algorithm to determine this unique $c$ value, and then the original message is given  by $m = \sqrt[3]{c}$.
    \end{enumerate}


% \subsection*{\Exercise}
%     \begin{enumerate}
%         \item Use similar strategies as in class to find all positive integers $x, y$ such that $$|12^{x} - 5^{y}| = 7.$$
%         {\small \textit{Hint: Consider congruencies modulo 4 and modulo 6}}.
%     \end{enumerate}

% \subsection*{Solution}
    % If $x, y$ satisfy $|12^{x} - 5^{y}| = 7$, then either $12^{x} - 5^{y} = 7$ or $12^{x} - 5^{y} = -7$
%     \begin{enumerate}
%         \item 
%         \begin{Cases}
%             \item  $5^{y} - 12^{x} = 7$. In modulo $4$, this becomes $1 \equiv 3 \pmod{4}$, so there are no solutions for this case.
%             \item $12^{x} - 5^{y} = 7$. In modulo $6$, this becomes $(-1)^{n} \equiv -1 \pmod{6}$, so $n$ is odd. Note that $m = n = 1$ is a solution, so taking modulo $5$ and $4$ would not help. Instead, we can assume that $m \ge 2$ and now we can analyze with modulo $8$: $$5^{n} \equiv 1 \pmod{8} \implies 2 \mid n.$$ This contradicts wtih $n$ begin odd, what we found earlier. So $m = n = 1$ is our only solution.
%         \end{Cases}
%     \end{enumerate}


% % Reused from last semester, not very elegant proof. Keep as backup.
% \subsection*{\Exercise}
%     Let $a, b \in \mathbb{N}^{*}$ with $a \ge b$.
%     \begin{enumerate}
%         \item Show that, if $a$ and $b$ are odd, $2\gcd(a, b) = \gcd(a + b, a - b)$.
%         \item Give a counterexample to show that the equality $2\gcd(a, b) = \gcd(a + b, a - b)$ is not always true.
%     \end{enumerate}

% \subsection*{Solution}
%     \begin{enumerate}
%         \item Let $d = 2\gcd(a, b)$. Then $\frac{d}{2} = \gcd(a, b)$, so $\frac{d}{2} \mid a$ and $\frac{d}{2} \mid b$, so $\frac{d}{2} \mid a + b = \frac{2(a + b)}{2}$ and $\frac{d}{2} \mid a - b = \frac{2(a - b)}{2}$. But $\frac{d}{2}$ and $2$ are coprime, so $\frac{d}{2} \mid \frac{a + b}{2}$ and $\frac{d}{2} \mid \frac{a - b}{2}$. Therefore $\frac{d}{2} \mid \gcd\left(\frac{a + b}{2}, \frac{a - b}{2}\right) = \frac{1}{2}\gcd(a + b, a - b)$, so $d \mid \gcd(a + b, a - b)$.
        
%         \vspace{1.5mm}
%         Now let $d' = \gcd(a + b, a - b)$. Then $d' \mid (a + b)$ and $d' \mid (a - b)$. Then $d' \mid 2a$ and $d' \mid 2b$, so $d' \mid \gcd(2a, 2b)$, so $d' \mid 2\gcd(a, b)$.
        
%         \vspace{1.5mm}
%         Now $d \mid d'$ and $d' \mid d$. Hence $d = d'$.
%         \item Let $a = 2$ and $b = 4$. Then $2\gcd(a, b) = 4$, but $\gcd(a + b, a - b) = 2$.
%     \end{enumerate}

% \subsection*{\Exercise \ (Not finished, will probably be removed, better suited for 3510)}
%     \textit{(Analysis of Algorithms)} Consider Euclid's algorithm given below, which calculates $\gcd(a, b)$ for two integers $a, b$ such that $0 \le b \le a$. % https://www.quora.com/What-is-the-time-complexity-of-Euclids-GCD-algorithm
%     \begin{enumerate}[leftmargin=1.2in]
%         \item [\textbf{Initialization.}] Set $a_{1} = a$ and $b_{1} = b$.
%         \item [\textbf{Inductive Step.}] Set $a_{i} = b_{i - 1}$, and set $b_{i} = a_{i - 1} \text{ (mod $b_{i - 1}$)}$. The cost of the $\text{mod}$ operation is constant $c$. Repeat this step, incrementing $i$ in each iteration, until $i = n$ such that $b_{n}$ is set equal to $0$.
%         \item [\textbf{Termination.}] Return $\gcd(a, b)$, given by $a_{n}$.
%     \end{enumerate}
%     In this exercise, we will find an upper bound on the worst-case time complexity of this algorithm.
%     \begin{enumerate}
%         \item Recall the Fibonacci sequence: $$f_{n} = f_{n - 1} + f_{n - 2}, f(1) = 0, f(2) = 1$$ Prove using induction that every two consecutive elements in this sequence are relatively prime. Use this fact to explain why the worst case for the Euclidean algorithm is on any two consecutive Fibonacci numbers.
%         % Prove using induction that the smallest values for $a_{i}$ and $b_{i}$ are $f_{n + 2}$ and $f_{n + 1}$ respectively.
%         \item It turns out the Fibonacci sequence has a closed form: $f_{i} = \frac{\phi^{i}}{\sqrt{5}}$. Use this information, along with part (a), to prove that an upper bound on the worst-case time complexity of this algorithm is $O(\log(b))$. Make sure to explain why the worst-case time complexity is dependent upon the \textit{smaller} of the two inputs.
%     \end{enumerate} 

% \subsection*{Solution}
%     \begin{enumerate}
%         \item Let $P(n)$ be the statement $$\gcd(f_{n + 1}, f_{n}) = 1.$$
%         \textsc{Base Case}: $P(1)$ is true, since $\gcd(f_{1}, f_{2}) = \gcd(1, 0) = 1$. \\
%         \textsc{Inductive Step}: Now let $n \in \mathbb{N}^{+}$ such that $P(n)$ is true. Using the inductive hypothesis, the Euclidean algorithm, and the fact that $f_{n + 2} = f_{n + 1} + f_{n}$, we obtain 
%         $$\gcd(f_{n + 2}, f_{n + 1}) = \gcd(f_{n + 1} + f_{n}, f_{n + 1}) = \gcd(f_{n + 1}, f_{n}) = 1.$$ Hence $P(n + 1)$ is true as well. \\
%         \textsc{Conclusion}: Therefore by induction, for all $n \in \mathbb{N}^{+}$, $P(n)$ is true.
        
%         \vspace{1.5mm}
%         \item 
%         \item Suppose at s fome iteration, we have $(a_{i}, b_{i})$. In the next iteration, we have $(b_{i}, a_{i} \text{ (mod $b_{i}$)})$. Note that after every iteration, both the numbers become at most the value of $b_{i}$.
        
        
%         More specifically, if $\gcd(a, b)$ requires n steps, and $a < b \leq 0$, then the smallest possible values of a and b are $f_{n+2}$ and $f_{n+1}$, respectively. This can be proven by induction.
        
%         So when $n$ iterations are required, we have that $b \ge f_{n + 1} \ge f_{n} = \frac{\phi^{n}}{\sqrt{5}}$. 
        
%         Finally, solving for $n$, we find that $n \le \log_{\phi}(b)$.
%     \end{enumerate}
    
% \subsection*{\Exercise}
%     (May be cut) Prove the following identity using the Binomial Theorem $$\sum_{i = 0}^{c}\binom{a}{c - i}\binom{b}{i} = \binom{a + b}{c}.$$  \textit{(Hint: calculate the coefficient of $x^{c}$ in $(x + 1)^{a}(x + 1)^{b} = (x + 1)^{a + b}$ in two different ways.)}

% \subsection*{Solution}
% By the binomial theorem, $$(1 + x)^{m + n} = \sum_{r=0}^{m+n} \binom{m + n}{r} x^{r}.$$
% Also by the binomial theorem,
% $$\begin{aligned}
%     (1 + x)^{m + n} &= (1 + x)^{m}(1 + x)^{n} \\
%     &= \left(\sum_{i=0}^{m} \binom{m}{i} x^{i}\right) \left(\sum_{j=0}^{n} \binom{n}{j} x^{j}\right) \\
%     &= \sum_{r = 0}^{m + n}\left(\sum_{k = 0}^{r} \binom{m}{k}\binom{n}{r - k}\right)x^{r} why? 
% \end{aligned}$$
% \textbf{Remark.} This is commonly known as \textit{Vandermonde's Identity}.


\end{document}